\section{Capitoli 20-21: Grafi}

\begin{dizionario}{Grafo Orientato vs Non Orientato}
\textbf{Orientato:} A $\to$ B (Senso unico).
\textbf{Non Orientato:} A $-$ B (Doppio senso).
\textbf{Ciclo:} Un percorso che parte da un nodo e torna allo stesso nodo.
\end{dizionario}

\begin{esercizio}{Es. 39: Grado di un Nodo}
\textbf{TRACCIA:} Dato un grafo (mappa di adiacenza) e un nodo, restituire il numero di vicini (grado).
\end{esercizio}

\begin{jsbox}
function grado(grafo, nodo) {
    if (!grafo.has(nodo)) return 0;
    return grafo.get(nodo).length;
}
\end{jsbox}

\begin{esercizio}{Es. 40: Esiste Percorso (DFS)}
\textbf{TRACCIA:} Scrivere `haPercorso(g, start, end)` che ritorna true se esiste una strada da `start` a `end`.
Usa DFS e un Set `visited` per evitare loop.
\end{esercizio}

\begin{jsbox}
function haPercorso(grafo, curr, target, visited = new Set()) {
    if (curr === target) return true;
    
    visited.add(curr);
    let vicini = grafo.get(curr) || [];
    
    for (let v of vicini) {
        if (!visited.has(v)) {
            if (haPercorso(grafo, v, target, visited)) return true;
        }
    }
    return false;
}
\end{jsbox}

\begin{esercizio}{Es. 41: Distanza Minima (BFS)}
\textbf{TRACCIA:} Trovare la distanza minima (numero di passi) tra `start` e `end`.
Usa BFS (Code).
\end{esercizio}

\begin{jsbox}
function distanzaMinima(grafo, start, end) {
    let queue = [[start, 0]]; // [nodo, distanza]
    let visited = new Set([start]);
    
    while (queue.length > 0) {
        let [curr, dist] = queue.shift();
        
        if (curr === end) return dist;
        
        let vicini = grafo.get(curr) || [];
        for (let v of vicini) {
            if (!visited.has(v)) {
                visited.add(v);
                queue.push([v, dist + 1]);
            }
        }
    }
    return -1; // irraggiungibile
}
\end{jsbox}
